% !TeX program = lualatex
% !TeX encoding = utf8
% !TeX spellcheck = uk_UA


% chktex-file 1
% chktex-file 8

%% --------------------------------------------------------
\Opensolutionfile{answer}[\currfilebase/\currfilebase-Answers]
\Writetofile{answer}{\protect\section*{\nameref*{\currfilebase}}}
\chapter{Система рівнянь Максвелла}\label{\currfilebase}\hypertarget{MaxwellEqns}{}
%% --------------------------------------------------------


\graphicspath{{\currfilebase/Pictures}}
%=========================================================
\epigraph{\Alexander Кількісне, математичне формулювання законів поля дано в~так званих
рівняннях Максвелла. Експериментальні факти призвели до формулювання цих рівнянь, але
зміст їх значно багатший [\ldots]. Їхня проста форма приховує глибину, що виявляється
тільки при ретельному вивченні. }{А. Ейнштейн, Л. Інфельд, \enquote{Еволюція фізики}}


%% ========================================================
\section*{Основоположники теорії електромагнітного поля}
%% ========================================================


Теорію електромагнітного поля, початки якої заклав \emph{Фарадей}, математично завершив
\emph{Максвелл} (рис.~\ref{pic:FaradayMaxwell}). Однією з~найважливіших нових ідей, висунутих Максвеллом, була думка про
\emph{симетрію} у~взаємозалежності електричного й~магнітного полів: якщо змінне магнітне
поле породжує електричне (закон Фарадея, розділ~\ref{ElMagInduction}), то, за красивою
симетрією, змінне електричне поле мусить породжувати магнітне. Ця ідея, як ми зараз
побачимо, узгоджується із законом збереження заряду й~веде до відкриття \emphx{струму
зміщення}~--- останньої відсутньої ланки повної системи рівнянь електродинаміки.

\begin{figure}[htbp]
	\centering
	\begin{minipage}[t]{0.45\linewidth}\centering
		\includegraphics[height=4.7cm]{Faraday}\\
		\footnotesize Майкл Фарадей\\ (1791--1867)
	\end{minipage}
	\hfill
	\begin{minipage}[t]{0.45\linewidth}\centering
		\includegraphics[height=4.7cm]{james-clerk-maxwell}\\
		\footnotesize Джеймс Клерк Максвелл\\ (1831--1879)
	\end{minipage}
	\caption{Основоположники класичної теорії електромагнітного
	поля\label{pic:FaradayMaxwell}}
\end{figure}

Формулювання рівнянь Максвелла А.~Ейнштейн і~Л.~Інфельд у~книзі \enquote{Еволюція фізики}
назвали \enquote{найважливішою подією з~часу Ньютона} --- не лише через цінність їхнього
змісту, а~й~тому, що вони дають зразок \emph{нового типу законів}: \emphx{рівняння
Максвелла~--- це закони, що виражають структуру поля}, а~не закони руху окремих частинок.


%% ========================================================
\section{Струм зміщення}\hypertarget{displacement-current}{}
%% ========================================================


% ---------------------------------------------------------
\subsection*{Суперечність у~законах магнетизму}
% ---------------------------------------------------------


Теорема про циркуляцію для постійного магнітного поля, яку ми встановили
в~розділі~\ref{MagFieldMedia},
% .........................................................
\begin{equation}\label{eq:CirculationHVacuumOld}
	\oint\limits_L \Hfield\cdot d\vect r = \frac{4\pi}{c}\iint\limits_S \vect j\cdot d\vect S, \qquad \rot\Hfield = \frac{4\pi}{c}\vect j,
\end{equation}
% .........................................................
виявляється \emph{неправильною} у~випадку \emph{змінного} електричного поля
(рис.~\ref{pic:CircThm}). Застосуємо операцію $\divg$ до диференціальної форми цього
рівняння: оскільки дивергенція будь-якого ротора тотожно дорівнює нулю, $\divg\rot\Hfield
\equiv 0$, отримуємо
% .........................................................
\begin{equation}\label{eq:DivJZero}
	\divg\vect j = 0.
\end{equation}
% .........................................................

\begin{SCfigure}[1][htbp]\centering
	\localinput{circ_thm.tikz}
	\caption{Контур $L$, натягнута на нього поверхня $S$ та провідник зі струмом, що її
	пронизує\label{pic:CircThm}}
\end{SCfigure}

Але, з~іншого боку, за законом збереження заряду (рівняння неперервності,
розділ~\ref{SteadyCurrent}), якщо густина заряду змінюється з~часом, $\partial\rho/\partial t
\neq 0$, то
% .........................................................
\begin{equation}\label{eq:ContinuityViolation}
	\divg\vect j = -\frac{\partial\rho}{\partial t} \neq 0.
\end{equation}
% .........................................................
Рівняння~\eqref{eq:DivJZero} і~\eqref{eq:ContinuityViolation} відверто \emph{суперечать}
одне одному щоразу, коли $\partial\rho/\partial t\neq0$ --- а~отже, теорему про
циркуляцію~\eqref{eq:CirculationHVacuumOld} необхідно \emph{видозмінити}.

Парадокс видно й~без дивергенції. Заряджаємо плоский конденсатор струмом $I$ (рис.~\ref{pic:TwoSurfaces}) і~беремо контур $L$, що охоплює дріт поблизу конденсатора. Натягнемо на нього дві поверхні: ліву, що перетинає дріт, і~праву, що проходить між обкладками й~не перетинає жодного дроту. За~\eqref{eq:CirculationHVacuumOld} циркуляція по лівій поверхні дорівнює $\frac{4\pi}{c}I$, а~по правій~--- нулю, бо струму провідності між обкладками немає. Але циркуляція по одному й~тому самому контуру не залежить від вибору поверхні (теорема Стокса): ми дістали суперечність.

\begin{figure}[htbp]\centering
	\localinput{two-surfaces-loop.tikz}
	\caption{Дві різні поверхні, натягнуті на один і~той самий контур $L$: ліва перетинається струмом провідності, права проходить між обкладками конденсатора\label{pic:TwoSurfaces}}
\end{figure}


% ---------------------------------------------------------
\subsection*{Гіпотеза Максвелла}
% ---------------------------------------------------------


Щоб\index{Максвелла гіпотеза} усунути цю суперечність, Дж.~Максвелл увів поняття \emphx{струму зміщення}\index{Струм зміщення} $\vect
j_\text{зм}$, доповнивши теорему про циркуляцію так, щоб закон збереження заряду
виконувався за будь-яких обставин:
% .........................................................
\begin{equation}\label{eq:DisplacementCurrentDef}
	\rot\Hfield = \frac{4\pi}{c}\left(\vect j+\vect j_\text{зм}\right).
\end{equation}
% .........................................................
Застосовуючи операцію $\divg$ до цього рівняння:
% .........................................................
\begin{equation*}
	\divg\left(\vect j+\vect j_\text{зм}\right) = 0 \quad\Rightarrow\quad \divg\vect j_\text{зм} = \frac{\partial\rho}{\partial t}
\end{equation*}
% .........................................................
(тут ми скористались рівнянням~\eqref{eq:ContinuityViolation}). За теоремою Гаусса для
електричного поля, $\rho = \dfrac1{4\pi}\divg\Dfield$ (розділ~\ref{FieldInDielectrics}), тому
% .........................................................
\begin{equation*}
	\divg\vect j_\text{зм} = \frac{\partial}{\partial t}\left(\frac1{4\pi}\divg\Dfield\right) = \divg\left(\frac1{4\pi}\frac{\partial\Dfield}{\partial
		t}\right).
\end{equation*}
% .........................................................
Найпростіше (і, як показує подальша перевірка експериментом, єдино правильне) розв'язання
цього рівняння~--- покласти
% .........................................................
\begin{equation}\label{eq:DisplacementCurrentDeriv}
	\tcbhighmath{
		\vect j_\text{зм} = \frac1{4\pi}\frac{\partial\Dfield}{\partial t}.
	}
\end{equation}
% .........................................................

Рівняння~$\divg\vect j_\text{зм} = \divg\left(\tfrac1{4\pi}\partial\Dfield/\partial
t\right)$ саме собою \emph{не доводить}, що $\vect j_\text{зм}$ дорівнює точно
$\tfrac1{4\pi}\partial\Dfield/\partial t$: до цього виразу можна було б, узагалі
кажучи, додати довільне бездивергентне поле ($\divg\vect f=0$). Максвелл обрав
\emph{найпростіший} можливий варіант~--- без жодного додаткового доданка; те, що саме
цей вибір відповідає дійсності (а~не просто математичній зручності),~--- це вже суто
\emph{експериментальний} факт, підтверджений, зокрема, дослідами Г.~Герца
з~випромінювання електромагнітних хвиль (див. розділ~\ref{ElMagWaves}).

\begin{Historical}
	Викладене вище виведення формули для струму зміщення~\eqref{eq:DisplacementCurrentDeriv} через закон збереження заряду~--- це \emph{логічна реконструкція}, зручна для навчального курсу, а~не той шлях, яким ішов сам Максвелл. Він прийшов до цієї ідеї через складну механічну модель ефіру (\enquote{молекулярні вихори} з~\enquote{холостими коліщатками} між ними), у~якій електричне поле спричиняло пружну деформацію~--- зміщення~--- середовища; звідси й~назва. Згодом від ефіру відмовилися, але назва \enquote{струм зміщення} лишилася за всією величиною $\vect j_\text{зм}$. Те, що рівняння можна отримати й~без механічних гіпотез, свідчить: вони відображають об'єктивну структуру поля, а~не модель, яка привела до їх відкриття.
\end{Historical}

Фізичний зміст $\vect j_\text{зм}$ видно з~означення $\Dfield=\Efield+4\pi\vect P$:
% .........................................................
\begin{equation*}
	\vect j_\text{зм} = \frac{1}{4\pi}\parttime{\Dfield} = \frac{1}{4\pi}\parttime{\Efield} + \parttime{\vect P}.
\end{equation*}
% .........................................................
Другий доданок, $\vect j_\text{пол} = \parttime{\vect P}$,~--- це \emphx{струм поляризації}, реальний рух зв'язаних зарядів у~речовині. Перший доданок описує швидкість зміни самого електричного поля й~існує навіть у~вакуумі.


% ---------------------------------------------------------
\subsection*{Теорема про циркуляцію магнітного поля з~урахуванням струму зміщення}
% ---------------------------------------------------------


Отже,\index{Теорема про циркуляцію!зі струмом зміщення} теорема про циркуляцію для магнітного поля, узгоджена із законом збереження заряду,
набуває остаточного вигляду:
% .........................................................
\begin{equation}\label{eq:CirculationHFull}
	\tcbhighmath{
		\rot\Hfield = \frac{4\pi}{c}\vect j + \frac1c\frac{\partial\Dfield}{\partial t}.
	}
\end{equation}
% .........................................................
В~інтегральній формі:
% .........................................................
\begin{equation}\label{eq:CirculationHFullInt}
	\tcbhighmath{
		\oint\limits_L \Hfield\cdot d\vect r = \frac{4\pi}{c}\iint\limits_S \vect j\cdot d\vect S + \frac1c\iint\limits_S
		\frac{\partial\Dfield}{\partial t}\cdot d\vect S,
	}
\end{equation}
% .........................................................
де
% .........................................................
\begin{equation}\label{eq:DisplacementCurrentIntegral}
	I_\text{зм} = \frac1{4\pi}\iint\limits_S \frac{\partial\Dfield}{\partial t}\cdot d\vect S
\end{equation}
% .........................................................
--- струм зміщення, що пронизує поверхню $S$, натягнуту на контур $L$. Тут $\vect j$~---
густина струму провідності (і~конвекційного струму, тобто руху заряджених тіл). Отже,
згідно з~гіпотезою Максвелла, \emph{змінне електричне поле поряд зі звичайними струмами
провідності також створює магнітне поле}.


% ---------------------------------------------------------
\subsection*{Порівняння закону Фарадея та гіпотези Максвелла}
% ---------------------------------------------------------


Порівняємо закон електромагнітної індукції Фарадея (розділ~\ref{ElMagInduction})
і~\enquote{оновлену} теорему про циркуляцію (закон Ампера--Максвелла) за відсутності струмів провідності ($\vect
j=0$) у~вакуумі ($\epsilon=\mu=1$, тобто $\Efield=\Dfield$, $\Hfield=\Bfield$):
% .........................................................
\begin{equation}\label{eq:FaradayVsMagnetoelectric}
	\underbrace{\rot\Efield = -\frac1c\frac{\partial\Bfield}{\partial t}}_\text{закон електромагнітної індукції} \qquad\qquad
	\underbrace{\rot\Bfield = +\frac1c\frac{\partial\Efield}{\partial t}}_\text{\enquote{закон магнітоелектричної індукції}}.
\end{equation}
% .........................................................

Ця пара рівнянь~--- вершина ідеї Максвелла про \emph{симетрію} електрики
й~магнетизму: \emph{змінне в~часі магнітне поле породжує вихрове електричне поле,
а~змінне в~часі електричне поле породжує магнітне поле}. Єдина відмінність~--- знак:
у~законі Фарадея \enquote{мінус} (правило Ленца), у~гіпотезі Максвелла~---
\enquote{плюс}. Ця відносно невелика асиметрія знаків, утім, має глибокі наслідки:
саме вона (разом з рівняннями $\divg\Efield=0$, $\divg\Bfield=0$ у~вакуумі) приводить,
як буде показано в~наступних розділах курсу, до існування \emphx{електромагнітних
хвиль}, що біжать зі швидкістю світла $c$, а~не до експоненційно загасаючих чи
зростаючих розв'язків.

\begin{figure}[htbp]\centering
	\localinput{radial-leak.tikz}
	\caption{Радіальне стікання заряду з~кулі: у~кожній точці густина струму провідності $\vect j$ і~густина струму зміщення $\vect j_\text{зм}$ рівні за модулем і~протилежні за напрямком, тож магнітного поля немає\label{pic:RadialLeak}}
\end{figure}

\begin{Example}{Радіальне стікання заряду з~кулі}
	Нехай куля несе заряд $Q(t)$, що стікає в~зовнішнє (слабко провідне) середовище
	рівномірно в~усіх напрямках (рис.~\ref{pic:RadialLeak}). Для простоти $\epsilon=1$. Чи виникає при цьому магнітне поле?

	Стікання заряду створює радіальний струм провідності $I = -\dfrac{\partial Q}{\partial
	t}$. Поле кулі на відстані $r$, за теоремою Гаусса, $\Efield(r) =
	\dfrac{Q(t)}{r^3}\vect r$ (і~теж спадає з~часом), тому крізь будь-яку сферу радіуса
	$r>R$ додатково тече струм зміщення:
	% .........................................................
	\begin{equation}\label{eq:RadialLeakCurrent}
		I_\text{зм} = \frac1{4\pi}\oiint\limits_S \frac{\partial\Efield}{\partial t}\cdot d\vect S = \frac{\partial Q}{\partial t} = -I.
	\end{equation}
	% .........................................................
	За теоремою про циркуляцію, застосованою до довільного кола $L$, що лежить на сфері
	радіуса $r>R$ (натягнемо на нього частину цієї ж сфери~--- поверхню $S$, яку
	перетинають і~струм провідності, і~струм зміщення):
	% .........................................................
	\begin{equation}\label{eq:RadialLeakZeroB}
		\oint\limits_L \Hfield\cdot d\vect r = \frac{4\pi}c\left(I+I_\text{зм}\right) = 0 \quad\Rightarrow\quad \Hfield=\Bfield=0.
	\end{equation}
	% .........................................................
	Чому з~нульової циркуляції випливає саме $\Hfield=0$? Коло на сфері довільне, тому дотичні до сфери компоненти поля дорівнюють нулю. Радіальна компонента через сферичну симетрію однакова в~усіх точках сфери. Водночас потік $\Bfield$ крізь замкнену сферу дорівнює нулю ($\divg\Bfield=0$, бо магнітних зарядів немає). Отже, і~радіальна компонента нульова.

	Так струм зміщення виявляється реальним внеском у~джерела магнітного поля, а~не поправкою \enquote{про око}: у~кожній точці він точно компенсує струм провідності.
\end{Example}




\begin{figure}[htbp]\centering
	\localinput{capacitor-charging-field.tikz}
	\caption{Магнітне поле, створене струмом зміщення всередині конденсатора під час його зарядки\label{pic:CapacitorChargingField}}
\end{figure}

\begin{Example}{Струм зміщення в~конденсаторі}
	Повернімося до конденсатора з~рис.~\ref{pic:TwoSurfaces}. Тепер для правої поверхні циркуляція визначається \emph{лише} струмом зміщення, $\oint_L\Hfield\cdot d\vect r = \tfrac{4\pi}c I_\text{зм}$, а~для лівої~--- струмом провідності, $\tfrac{4\pi}{c}I$. Циркуляція по одному контуру не залежить від вибору поверхні, тому
	% .........................................................
	\begin{equation}\label{eq:CapacitorTwoSurfaces}
		I = I_\text{зм},
	\end{equation}
	% .........................................................
	тобто між обкладками \enquote{протікає} струм зміщення, який замикає коло. Розглянемо докладніше, яке магнітне поле він створює (рис.~\ref{pic:CapacitorChargingField}).

	Нехай плоский конденсатор із пластинами радіуса $R$ заряджається струмом $I$ (зазор
	між пластинами малий, крайовими ефектами нехтуємо). Густина струму зміщення між
	обкладками дорівнює
	% .........................................................
	\begin{equation*}
		j_\text{зм} = \frac{I}{\pi R^2},
	\end{equation*}
	% .........................................................
	і~він спрямований так само як струм провідності в~підвідних дротах. Магнітне поле
	всередині конденсатора знаходимо за теоремою про циркуляцію, взявши контур радіуса $r$
	із центром на осі конденсатора:
	% .........................................................
	\begin{equation*}
		\oint_L \Hfield\cdot d\vect r = \frac{4\pi}{c}I_\text{зм}(r),
	\end{equation*}
	% .........................................................
	де $I_\text{зм}(r) = j_\text{зм}\pi r^2 = I r^2/R^2$. Звідси
	% .........................................................
	\begin{equation*}
		2\pi r H(r) = \frac{4\pi}{c}I\frac{r^2}{R^2}
		\quad\Rightarrow\quad
		H(r) = \frac{2Ir}{cR^2}.
	\end{equation*}
	% .........................................................

	Отже, всередині конденсатора виникає вихрове магнітне поле, силові лінії якого
	охоплюють вісь конденсатора за правилом правого гвинта. Це поле точнісінько таке саме,
	як у~провіднику того самого радіуса зі струмом $I$: у~цьому й~полягає фізичний зміст
	твердження, що струм зміщення \emph{замикає} коло й~створює магнітне поле нарівні зі
	струмом провідності.

	\emph{Числова оцінка.} Для струму зарядки $I=1$~А і~пластин радіуса $R=5$~см поле біля краю
	пластин $B(R)=2I/(cR)\approx0{,}04$~Гс ($4$~мкТл), що приблизно в~десять разів слабше за магнітне
	поле Землі. Саме тому Максвелл не міг перевірити свою гіпотезу безпосередньо: таке поле
	вперше виміряли лише 1929~року, а~переконливим підтвердженням стали електромагнітні хвилі
	Герца (1888).
\end{Example}


\emphx{Струми зміщення замикають струми провідності і~створюють магнітне поле
точнісінько так само, як і~звичайні струми провідності}
(рис.~\ref{pic:CapacitorChargingField}). Проте вони мають і~принципову відмінність:
струм зміщення (у~вакуумі) \emph{не} пов'язаний з~тепловиділенням~--- до нього
незастосовні ні закон Ома, ні закон Джоуля--Ленца (розділ~\ref{SteadyCurrent}), адже він
не пов'язаний з~переміщенням реальних носіїв заряду проти сил опору, а~лише зі зміною
електричного поля в~просторі. (У~реальному діелектрику змінна поляризація
супроводжується діелектричними втратами, але це властивість речовини, а~не струму
зміщення як такого.)




%% ========================================================
\section{Система рівнянь Максвелла}\hypertarget{maxwell-equations-system}{}
%% ========================================================


Доповнивши\index{Максвелла рівняння} основні закони електромагнетизму, встановлені в~попередніх розділах, гіпотезою
про струм зміщення, Максвелл отримав повну систему фундаментальних рівнянь
електродинаміки. Таких рівнянь~--- чотири. У~диференціальній формі
% .........................................................
\begin{subequations}\label{eq:MaxwellDiff}
	\begin{align}
		\divg\Dfield &= 4\pi\rho,\label{eq:MaxwellGaussE}\\
		\divg\Bfield &= 0,\label{eq:MaxwellGaussB}\\
		\rot\Efield &= -\frac1c\frac{\partial\Bfield}{\partial t},\label{eq:MaxwellFaraday}\\
		\rot\Hfield &= \frac{4\pi}{c}\vect j + \frac1c\frac{\partial\Dfield}{\partial t},\label{eq:MaxwellAmpere}
	\end{align}
\end{subequations}
% .........................................................
а~в~інтегральній (з~неї виводять граничні умови)
% .........................................................
\begin{align}
	\oiint\limits_S \Dfield\cdot d\vect S &= 4\pi\iiint\limits_V \rho\,dV,\label{eq:MaxwellGaussEInt}\\
	\oiint\limits_S \Bfield\cdot d\vect S &= 0,\label{eq:MaxwellGaussBInt}\\
	\oint\limits_L \Efield\cdot d\vect r &= -\frac1c\iint\limits_S \frac{\partial\Bfield}{\partial t}\cdot d\vect S,\label{eq:MaxwellFaradayInt}\\
	\oint\limits_L \Hfield\cdot d\vect r &= \frac{4\pi}{c}\iint\limits_S\left(\vect j + \frac1{4\pi}\frac{\partial\Dfield}{\partial t}\right)\cdot d\vect S.\label{eq:MaxwellAmpereInt}
\end{align}
% .........................................................
Фізичний зміст кожного рівняння такий.
\begin{itemize}
	\item \emph{Теорема Гаусса для $\Dfield$}~\eqref{eq:MaxwellGaussE}. Джерела електричного поля~--- вільні заряди: лінії $\Dfield$ починаються й~закінчуються на них, а~потік крізь замкнену поверхню дорівнює $4\pi$, помноженому на заряд усередині.
	\item \emph{Теорема Гаусса для $\Bfield$}~\eqref{eq:MaxwellGaussB}. Магнітних зарядів немає: лінії $\Bfield$ замкнені, і~потік крізь будь-яку замкнену поверхню дорівнює нулю.
	\item \emph{Закон індукції Фарадея}~\eqref{eq:MaxwellFaraday}. Змінне магнітне поле породжує вихрове електричне: циркуляція $\Efield$ дорівнює швидкості зменшення потоку $\Bfield$ (з~множником $1/c$).
	\item \emph{Теорема про циркуляцію зі струмом зміщення}~\eqref{eq:MaxwellAmpere}. Магнітне поле породжують і~струми провідності, і~змінне електричне поле.
\end{itemize}
Рисунок~\ref{pic:MaxwellScheme} підсумовує цю структуру: заряди й~струми~--- зовнішні джерела, а~змінні $\Efield$ і~$\Bfield$ породжують одне одного.

\begin{figure}[htbp]\centering
	\localinput{maxwell-scheme.tikz}
	\caption{Структура рівнянь Максвелла: заряди й~струми породжують поля, а~змінні електричне й~магнітне поля породжують одне одного; магнітних зарядів немає\label{pic:MaxwellScheme}}
\end{figure}

Тут $\rho$ і~$\vect j$~--- густини вільних зарядів і~струмів (зв'язані заряди й~струми
враховано через $\Dfield$ і~$\Hfield$). Закон збереження заряду не є окремим рівнянням,
а~випливає з~\eqref{eq:MaxwellGaussE} і~\eqref{eq:MaxwellAmpere}: застосувавши $\divg$ до~\eqref{eq:MaxwellAmpere} й~підставивши~\eqref{eq:MaxwellGaussE}, отримуємо $\divg\vect j + \partial\rho/\partial t = 0$.

\begin{Historical}
	У~такому компактному вигляді рівняння Максвелла \emph{самому Максвеллу невідомі}.
	У~статті \enquote{Динамічна теорія електромагнітного поля} (1865) теорія викладена
	двадцятьма рівняннями для двадцяти змінних, записаними покомпонентно, а~в~\enquote{Трактаті про електрику
	й~магнетизм} (1873)~--- частково в~кватерніонній нотації. Чотири векторні рівняння вигляду, наведеного
	в~системі~\eqref{eq:MaxwellDiff}, виокремили Олівер Хевісайд (1884--1885~рр.) і~незалежно від нього
	Генріх Герц (1890~р.), а~векторний апарат $\Grad$, $\Div$, $\Rot$ розвивали водночас із ними
	Дж.~В.~Гіббс і~сам Хевісайд. Тому їх іноді називають \enquote{рівняннями
	Герца--Хевісайда}.
\end{Historical}


% ---------------------------------------------------------
\subsection*{Вираз електричного поля через потенціали}
% ---------------------------------------------------------


Рівняння Максвелла природно групуються парами: \emph{перша} пара (теорема Гаусса для
$\Bfield$ і~циркуляція $\Efield$)~--- рівняння \emph{без} явних джерел (зарядів
і~струмів), \emph{друга} пара~--- рівняння \emph{з} джерелами.

Теорема Гаусса для $\Bfield$ дає змогу, як і~в~магнітостатиці (розділ~\ref{MagFieldVacuum}),
ввести вектор-потенціал:
% .........................................................
\begin{equation*}
	\Bfield = \rot\vect A.
\end{equation*}
% .........................................................
Підставляючи це в~закон індукції, $\rot\Efield = -\dfrac1c\parttime{\Bfield} =
-\dfrac1c\parttime{}(\rot\vect A) = -\rot\left(\dfrac1c \parttime{\vect A}\right)$,
отримуємо
% .........................................................
\begin{equation*}
	\rot\left(\Efield+\frac1c\parttime{\vect A}\right) = 0.
\end{equation*}
% .........................................................
Оскільки ротор цього виразу тотожно дорівнює нулю, сам вираз можна подати як градієнт
деякої скалярної функції $-\phi$:
% .........................................................
\begin{equation}\label{eq:EFromPotentials}
	\tcbhighmath{
		\Efield = -\grad\phi - \frac1c\frac{\partial\vect A}{\partial t}.
	}
\end{equation}
% .........................................................
У~випадку сталих у~часі полів $\parttime{\vect A}=0$, і~$\Efield=-\grad\phi$~--- введена
тут функція $\phi$ природно збігається зі звичним скалярним потенціалом електростатики.
Зауважимо, що потенціали визначені неоднозначно: заміна $\vect A\to\vect A+\grad f$,
$\phi\to\phi-\frac1c\parttime{f}$ не змінює ні $\Bfield$, ні $\Efield$ (калібрувальна
свобода, розділ~\ref{MagFieldVacuum}).

Формула~\eqref{eq:EFromPotentials} показує, що в~нестаціонарному випадку електричне поле визначається \emph{не тільки} скалярним потенціалом $\phi$, а~й швидкістю зміни $\vect A$: ця частина $\Efield$ породжується не зарядами напряму, а~змінним магнітним полем. Це й~є вихрова, \enquote{індукційна} частина поля (у~кулонівському калібруванні $\divg\vect A = 0$), з~якою ми зустрічались у~розділі~\ref{ElMagInduction}.


% ---------------------------------------------------------
\subsection*{Стаціонарні поля}
% ---------------------------------------------------------


У~стаціонарному випадку всі часткові похідні за часом дорівнюють нулю, $\parttime{}=0$,
і~повна система рівнянь Максвелла розпадається на дві \emph{незалежні} системи~--- для
електростатики й~магнітостатики (рис.~\ref{pic:StaticsComparison}):
% .........................................................
\begin{equation}\label{eq:StaticDecoupling}
	\underbrace{\begin{cases}
			\divg\Dfield = 4\pi\rho, \\
			\rot\Efield = 0.
		\end{cases}}_\text{електростатика} \qquad\qquad
	\underbrace{\begin{cases}
			\divg\Bfield = 0, \\
			\rot\Hfield = \dfrac{4\pi}c\vect j.
		\end{cases}}_\text{магнітостатика}
\end{equation}
% .........................................................

\begin{figure}[htbp]\centering
	\begin{subfigure}[t]{0.48\linewidth}\centering
		\localinput{elstat.tikz}
		\caption{Електростатика: поле точкового заряду}
		\label{tikz:elstat}
	\end{subfigure}\hfill
	\begin{subfigure}[t]{0.48\linewidth}\centering
		\localinput{magstat.tikz}\\
		\caption{Магнітостатика: поле провідника зі струмом}
		\label{tikz:magstat}
	\end{subfigure}
	\caption{У~стаціонарному випадку електричне (\protect\subref{tikz:elstat}) й~магнітне
	(\protect\subref{tikz:magstat}) поля \enquote{не знають} одне про
	одного\label{pic:StaticsComparison}}
\end{figure}

Отже, усі попередні розділи курсу~--- електростатика й~магнітостатика~--- не окремі теорії, а~\emph{частинні випадки} єдиної теорії Максвелла для полів, що не залежать від часу. \emph{Зв'язок} електрики й~магнетизму (взаємне породження) принципово \emph{нестаціонарний}: він зникає, щойно всі процеси застигають у~часі.


% ---------------------------------------------------------
\subsection*{Матеріальні рівняння}
% ---------------------------------------------------------


Рівняння Максвелла~--- це рівняння \emph{загальної} структури поля: чотири векторні рівняння містять п'ять векторів $\Efield$, $\Dfield$, $\Bfield$, $\Hfield$, $\vect j$, тож для розв'язання конкретної задачі їх потрібно доповнити \emphx{матеріальними рівняннями}\index{Матеріальні рівняння}, які пов'язують ці поля й~залежать від властивостей середовища. Допоміжні вектори $\Dfield$ і~$\Hfield$ означено через поляризацію $\vect P$ й~намагніченість $\vect J$ (розділи~\ref{FieldInDielectrics} і~\ref{MagFieldMedia}):
% .........................................................
\begin{equation}\label{eq:DHDefinition}
		\Dfield = \Efield + 4\pi\vect P, \qquad \Hfield = \Bfield - 4\pi\vect J.
\end{equation}
% .........................................................
Для лінійних ізотропних середовищ $\vect P=\chi\Efield$, $\vect J=\chi_m\Hfield$, і~система замикається співвідношеннями
% .........................................................
\begin{equation}\label{eq:ConstitutiveRelations}
		\Dfield = \epsilon\Efield, \qquad \Bfield = \mu\Hfield, \qquad \vect j = \lambda\Efield,
\end{equation}
% .........................................................
де $\epsilon=1+4\pi\chi$ і~$\mu=1+4\pi\chi_m$. Кожне з~них має свій фізичний зміст: $\epsilon$ описує відгук зв'язаних зарядів діелектрика (поляризацію), $\mu$~--- відгук елементарних струмів магнетика (намагнічення), а~$\lambda$~--- відгук вільних носіїв (закон Ома в~диференціальній формі, розділ~\ref{SteadyCurrent}). У~вакуумі $\epsilon=\mu=1$, $\lambda=0$. Для анізотропних середовищ $\epsilon$ і~$\mu$ стають тензорами, у~сильних полях залежать від поля, а~за швидкої зміни поля~--- від частоти, $\epsilon=\epsilon(\omega)$.


% ---------------------------------------------------------
\subsection*{Граничні умови}
% ---------------------------------------------------------


Диференціальні рівняння Максвелла описують поле всередині однорідних областей. На межі розділу двох середовищ $\epsilon$, $\mu$ (а~часом і~саме поле) змінюються стрибком, тож рівняння доповнюють \emphx{граничними умовами}\index{Граничні умови!для електромагнітного поля}. Їх виводять з~\emph{інтегральної} форми~\eqref{eq:MaxwellInt} так само, як в~електростатиці (\S~\ref{sec:BC_vacuum}) і~магнітостатиці (розділ~\ref{MagFieldMedia}): беруть нескінченно малий циліндр і~нескінченно вузький контур на межі. Єдине нове питання~--- чи не змінює змінність полів, $\partial\Bfield/\partial t$ і~$\partial\Dfield/\partial t$, вигляду цих умов. Нормаль $\vect n$ напрямлена із середовища~1 у~середовище~2.


% ---------------------------------------------------------
\subsection*{Нормальні складові}
% ---------------------------------------------------------


Застосуємо теореми Гаусса~\eqref{eq:MaxwellGaussEInt}, \eqref{eq:MaxwellGaussBInt} до нескінченно малого циліндра, що охоплює ділянку межі (рис.~\ref{pic:BCPillbox}). Висоту циліндра спрямуємо до нуля: потоком крізь бічну поверхню можна знехтувати, лишаються потоки крізь основи площею $dS$. Оскільки $d\vect S_2=\vect n\,dS$, $d\vect S_1=-d\vect S_2$, а~вільний заряд усередині циліндра $q=\sigma\,dS$ ($\sigma$~--- поверхнева густина вільних зарядів),
% .........................................................
\begin{align*}
	\oiint\limits_S \Dfield\cdot d\vect S = 4\pi q \quad&\Rightarrow\quad (\Dfield_2-\Dfield_1)\cdot\vect n\,dS = 4\pi\sigma\,dS,\\
	\oiint\limits_S \Bfield\cdot d\vect S = 0 \quad&\Rightarrow\quad (\Bfield_2-\Bfield_1)\cdot\vect n\,dS = 0.
\end{align*}
% .........................................................

\begin{figure}[htbp]\centering
	\localinput{bc-pillbox.tikz}
	\caption{Циліндр Гаусса на межі розділу двох середовищ: $\sigma$~--- вільні заряди на межі, $d\vect S_2=-d\vect S_1=\vect n\,dS$\label{pic:BCPillbox}}
\end{figure}

Отже,
% .........................................................
\begin{equation}\label{eq:BCNormal}
	\tcbhighmath{
		D_{2n} - D_{1n} = 4\pi\sigma, \qquad B_{2n} = B_{1n}.
	}
\end{equation}
% .........................................................
Нормальна складова $\Dfield$ стрибає на величину $4\pi\sigma$, якщо на межі є вільні заряди, і~неперервна без них. Нормальна складова $\Bfield$ неперервна завжди: магнітних зарядів немає. Рівняння Гаусса не містять похідних за часом, тож змінність полів цих умов не змінює.


% ---------------------------------------------------------
\subsection*{Тангенціальні складові}
% ---------------------------------------------------------


Застосуємо рівняння для циркуляції~\eqref{eq:MaxwellFaradayInt} і~\eqref{eq:MaxwellAmpereInt} до вузького прямокутного контуру завдовжки $d\ell$. Його довгі сторони проходять по обидва боки від межі, а~висота $h\to0$ (рис.~\ref{pic:BCLoop}). Нехай $\vect\tau$~--- одиничний вектор, дотичний до межі й~напрямлений уздовж верхньої сторони, тож $d\vect r_2=\vect\tau\,d\ell$, а~$d\vect r_1=-d\vect r_2$. Внеском коротких сторін нехтуємо. Для $\Efield$
% .........................................................
\begin{align*}
	\oint\limits_L \Efield\cdot d\vect r &= (\Efield_2-\Efield_1)\cdot\vect\tau\,d\ell\\
	&= -\frac1c\iint\limits_S \frac{\partial\Bfield}{\partial t}\cdot d\vect S,
\end{align*}
% .........................................................
і~права частина прямує до нуля: потік скінченної величини $\partial\Bfield/\partial t$ крізь площу $h\,d\ell\to0$ зникає разом з~нею. Тому тангенціальна складова $\Efield$ неперервна, і~змінність $\Bfield$ цього не змінює.

\begin{figure}[htbp]\centering
	\localinput{bc-loop.tikz}
	\caption{Контур циркуляції на межі розділу: поверхневий струм з~лінійною густиною $\vect i$ пронизує контур незалежно від його висоти $h$\label{pic:BCLoop}}
\end{figure}

Для $\Hfield$ праворуч стоїть ще й~струм: потік струму зміщення $\partial\Dfield/\partial t$ крізь площу $h\,d\ell$ теж зникає, а~струм провідності зникає, лише якщо він розподілений в~об'ємі. Якщо ж він тече тонким шаром на самій межі, то його \emphx{лінійна густина} $\vect i$ (струм через одиницю довжини, перпендикулярної до нього) від $h$ не залежить і~пронизує контур. Орієнтуємо контур упоперек струму, $\vect\tau=\vect b\times\vect n$, де $\vect b=\vect i/i$:
% .........................................................
\begin{equation*}
	\oint\limits_L \Hfield\cdot d\vect r = (\Hfield_2-\Hfield_1)\cdot\vect\tau\,d\ell = \frac{4\pi}{c}\,i\,d\ell.
\end{equation*}
% .........................................................
Оскільки $(\Hfield_2-\Hfield_1)\cdot(\vect b\times\vect n)=\vect b\cdot[\vect n\times(\Hfield_2-\Hfield_1)]$, а~контур іншої орієнтації струм не пронизує,
% .........................................................
\begin{equation}\label{eq:BCTangential}
	\tcbhighmath{
		E_{2\tau} = E_{1\tau}, \qquad \vect n\times\left(\Hfield_2 - \Hfield_1\right) = \frac{4\pi}{c}\,\vect i.
	}
\end{equation}
% .........................................................
Тангенціальна складова $\Efield$ неперервна завжди, а~тангенціальна складова $\Hfield$ стрибає лише там, де по межі течуть поверхневі струми провідності (наприклад, обмотка на межі двох магнетиків або, в~ідеалізації $\delta\to0$, скін-шар ідеального провідника, див. нижче); без них $H_{2\tau}=H_{1\tau}$.

Підсумуємо фізичний зміст умов~\eqref{eq:BCNormal} і~\eqref{eq:BCTangential}: стрибок $D_n$ створюють поверхневі вільні заряди, стрибок $H_\tau$~--- поверхневі струми; $B_n$ і~$E_\tau$ неперервні завжди. Для змінних полів умови мають \emph{той самий вигляд}, що й~для стаціонарних, бо струм зміщення й~$\partial\Bfield/\partial t$ скінченні й~у~тонкому шарі нічого не дають.


%% ========================================================
\section{Про рівняння Максвелла}\hypertarget{about-maxwell-equations}{}
%% ========================================================


\emph{Теорія Максвелла~--- це вершина класичної науки про електромагнітне поле}. До цієї
висоти ми піднімалися крок за кроком, уточнюючи поняття й~розкриваючи глибинні принципи
цієї науки: кожен пройдений розділ~--- електростатика, закони постійного струму,
магнітостатика, явище електромагнітної індукції~--- був немов базовий табір на шляху до
вершини (рис.~\ref{pic:Mountain}).

\begin{figure}[htbp]\centering
	\includegraphics[width=\linewidth]{Mountain}
	\caption{Теорія Максвелла~--- це вершина класичної науки про електромагнітне поле}
	\label{pic:Mountain}
\end{figure}

На вершині перед нами постають рівняння Максвелла~--- велична симфонія електричного
й~магнітного полів, ключ до розуміння електромагнітної взаємодії. Уся споруда класичної
електродинаміки, зведена нами розділ за розділом~--- від закону Кулона до струму
зміщення,~--- вміщується тепер у~чотирьох компактних рівняннях Максвелла, доповнених
трьома матеріальними рівняннями середовища. Ці чотири рівняння містять \emph{усе}:
електростатику й~магнітостатику як окремі стаціонарні випадки, електромагнітну індукцію
й~струм зміщення як дві половини єдиної симетрії,~--- і, як буде показано далі,
\emphx{електромагнітні хвилі} й~саме \emphx{світло} як прямий наслідок цих же рівнянь
у~вакуумі. Людвіг Больцман, автор одного з~перших викладів теорії Максвелла в~курсі лекцій (\enquote{Лекції з~теорії Максвелла про електрику й~світло}, 1891--1893), поставив епіграфом до другої частини слова Гете з~\enquote{Фауста}:

\begin{tornpage}\small\itshape
	Чи не бог накреслив ці знаки\ldots{}?

	\hfill{Й.~В.~Гете, \enquote{Фауст} (переклад наш)}.
\end{tornpage}

\vspace*{1.5ex}

Це не просто захоплення: вже за життя автора з~рівнянь випливло значно більше, ніж у~них було закладено. Максвелл записав систему, спираючись на закони Кулона, Ампера, Фарадея~--- і~жодне з~цих окремих досягнень електростатики та магнітостатики не містило натяку на те, що об'єднана система виявиться незрівнянно багатшою за суму своїх частин. Наслідки рівнянь подібні до спуску з~вершини у~простір застосувань. Перші кроки цього спуску ми зробимо вже в~межах цього розділу: закон
збереження енергії електромагнітного поля (теорема Пойнтінга) і~скін-ефект
у~провідниках~--- найближчі до вершини застосування щойно записаної системи. Далі,
у~наступних розділах курсу, цей спуск продовжиться~--- до поширення електромагнітних
хвиль, принципів їх випромінювання і, зрештою, до основ оптики, радіо- та електротехніки.


%% ========================================================
\section{Енергія електромагнітного поля}\label{sec:PintingTheoreme}\hypertarget{em-field-energy}{}
%% ========================================================


% ---------------------------------------------------------
\subsection*{Закон збереження енергії електромагнітного поля}
% ---------------------------------------------------------


Нехай\index{Енергія електромагнітного поля!закон збереження} у~деякій області простору присутнє електромагнітне поле й~заряджені частинки, що
з~ним взаємодіють (рис.~\ref{pic:EnergyFlowRegion}). Якщо ця система \emph{замкнена}, її
енергія зберігається; якщо ж \emph{незамкнена}, енергія може як \emph{втікати} в~цю
область, так і \emph{витікати} з~неї. Якісно закон збереження енергії для такої області
можна записати так:
% ========================================================
\begin{multline*}
	(\text{зменшення енергії поля в об'ємі}) = \\  (\text{робота над частинками}) + (\text{потік енергії назовні}),
\end{multline*}
% ========================================================
або, виражаючи цікаву для нас величину~--- роботу поля над зарядженими частинками:
% ========================================================
\begin{multline*}
	(\text{робота над частинками}) =  \\(\text{зменшення енергії поля в об'ємі}) - (\text{потік енергії назовні}).
\end{multline*}
% ========================================================

\begin{SCfigure}[2][htbp]\centering
	\localinput{energy-flow-region.tikz}
	\caption{Заряджені частинки в~області, обмеженій замкненою поверхнею: енергія
	електромагнітного поля може втікати й~витікати крізь цю
	поверхню\label{pic:EnergyFlowRegion}}
\end{SCfigure}

Рухати заряджені частинки, як ми твердо встановили ще в~розділі~\ref{MagFieldVacuum}, може
\emph{лише електричне} поле (магнітне поле роботи не виконує). Енергія, яку електричне
поле витрачає на розгін частинок в~одиниці об'єму за одиницю часу,~--- це потужність
$\vect j\cdot\Efield$ (у~провіднику вона перетворюється на джоулеве тепло,
розділ~\ref{SteadyCurrent}).


% ---------------------------------------------------------
\subsection*{Виведення закону збереження}
% ---------------------------------------------------------


Запишемо закон Максвелла для $\rot\Hfield$~\eqref{eq:CirculationHFull}, розв'язавши його
відносно $\vect j$:
% .........................................................
\begin{equation*}
	\vect j = \frac{c}{4\pi}\rot\Hfield - \frac1{4\pi}\parttime{\Dfield}.
\end{equation*}
% .........................................................
Помножимо скалярно обидві частини на $\Efield$:
% .........................................................
\begin{equation*}
	\vect j\cdot\Efield = \frac{c}{4\pi}\Efield\cdot\rot\Hfield - \frac1{4\pi}\Efield\cdot\parttime{\Dfield}.
\end{equation*}
% .........................................................
Скористаємось тотожністю векторного аналізу $\Efield\cdot\rot\Hfield =
-\divg\left[\Efield\times\Hfield\right] + \Hfield\cdot\rot\Efield$ та матеріальним
рівнянням $\Dfield=\epsilon\Efield$ (середовище лінійне, без дисперсії, $\epsilon$ і~$\mu$
сталі; тому $\Efield\cdot\parttime{\Dfield} = \epsilon\Efield\cdot\parttime{\Efield} =
\parttime{}\left(\tfrac12\epsilon E^2\right)$):
% .........................................................
\begin{equation*}
	\vect j\cdot\Efield = -\divg\left(\frac{c}{4\pi}\left[\Efield\times\Hfield\right]\right) + \frac{c}{4\pi}\Hfield\cdot\rot\Efield -
	\parttime{}\left(\frac1{8\pi}\epsilon E^2\right).
\end{equation*}
% .........................................................
Тепер підставимо закон електромагнітної індукції, $\rot\Efield =
-\dfrac1c\parttime{\Bfield}$, і~матеріальне рівняння $\Bfield=\mu\Hfield$ (так що
$\Hfield\cdot\parttime{\Bfield} = \parttime{}\left(\tfrac12\mu H^2\right)$):
% .........................................................
\begin{equation*}
	\vect j\cdot\Efield = -\divg\left(\frac{c}{4\pi}\left[\Efield\times\Hfield\right]\right) - \parttime{}\left(\frac{\epsilon E^2}{8\pi} +
	\frac{\mu H^2}{8\pi}\right).
\end{equation*}
% .........................................................


% ---------------------------------------------------------
\subsection*{Теорема Пойнтінга}
% ---------------------------------------------------------


Проаналізуємо отриманий вираз. Позначимо
% .........................................................
\begin{equation}\label{eq:PoyntingVectorDef}
	\tcbhighmath{
		\vect\Pi = \frac{c}{4\pi}\left[\Efield\times\Hfield\right]
	}
\end{equation}
% .........................................................
--- \emphx{вектор Пойнтінга}\index{Пойнтінга вектор}, і
% .........................................................
\begin{equation}\label{eq:energy_density}
	\tcbhighmath{
		w = \dfrac{\epsilon E^2}{8\pi}+\dfrac{\mu H^2}{8\pi} = \dfrac{\Efield\cdot\Dfield}{8\pi}+\dfrac{\Bfield\cdot\Hfield}{8\pi}
	}
\end{equation}
% .........................................................
--- об'ємну густину енергії електромагнітного поля (сума вже відомих нам електричної
й~магнітної частин, розділи~\ref{ElEnergy} і~\ref{MagEnergy}). Тоді
% .........................................................
\begin{equation}\label{eq:PoyntingTheoremDiff}
	\tcbhighmath{
		-\parttime{w} = \divg\vect\Pi + \vect j\cdot\Efield.
	}
\end{equation}
% .........................................................
Інтегруючи по довільному об'єму $V$, обмеженому поверхнею $S$, і~застосовуючи теорему
Остроградського--Гаусса, отримуємо \emph{інтегральну} форму~--- \emphx{теорему Пойнтінга}\index{Пойнтінга теорема}:
% .........................................................
\begin{equation}\label{eq:PoyntingTheoremInt}
	\tcbhighmath{
		-\frac{dW}{dt} = \oiint\limits_S \vect\Pi\cdot d\vect S + \iiint\limits_V \vect j\cdot\Efield\, dV.
	}
\end{equation}
% .........................................................

\emph{Теорема Пойнтінга}: зменшення енергії електромагнітного поля за одиницю часу в~даному об'ємі дорівнює сумі потоку енергії крізь поверхню, що обмежує цей об'єм, і~потужності, яку сили поля виконують над зарядами речовини всередині об'єму. Вектор $\vect\Pi$ визначає кількість енергії поля, що протікає крізь одиничну площадку за одиницю часу, і~вказує \emph{напрям} руху енергії~--- поняття, потрібне лише для змінних полів (у~статиці енергія нікуди не тече). Сам $\vect\Pi$ визначений з~точністю до ротора довільного векторного поля, але потік крізь замкнену поверхню однозначний.

\begin{Historical}
	Джон Генрі Пойнтінг (1852--1914), учень Максвелла, а~згодом перший професор фізики
	в~Бірмінгемі, опублікував теорему й~вектор, що носять його ім'я, 1884~р. у~статті
	\enquote{Про передачу енергії в~електромагнітному полі}. Головна нова думка~---
	енергія передається \emph{не вздовж провідника}, а~\emph{через простір навколо
	нього}: саме до цього висновку ми прийдемо нижче в~прикладах з~провідником
	і~конденсатором. Приблизно в~той самий час до такого ж результату незалежно дійшов
	Олівер Хевісайд, тож у~літературі трапляється й~назва \enquote{вектор Умова--Пойнтінга}
	(М.~О.~Умов ще 1874~р. запровадив поняття потоку енергії для пружних середовищ).
\end{Historical}


% ---------------------------------------------------------
\subsection*{Про одиниці вимірювання}
% ---------------------------------------------------------


У~гауссовій системі одиниць
% .........................................................
\begin{equation*}
	\left[\Pi\right] = \frac{\text{ерг}}{\text{см}^2\cdot\text{с}}, \qquad \left[w\right] = \frac{\text{ерг}}{\text{см}^3};
\end{equation*}
% .........................................................
для порівняння, у~СІ вектор Пойнтінга вимірюється у~Вт/м$^2$, а~густина енергії~---
у~Дж/м$^3$.

\begin{figure}[htbp]\centering
	\localinput{capacitor-poynting.tikz}
	\caption{Поля $\Efield$, $\Hfield$ і~вектор Пойнтінга $\vect\Pi$ при зарядці плоского конденсатора\label{pic:CapacitorPoynting}}
\end{figure}

\begin{Example}{Зарядка конденсатора}
	Плоский конденсатор з~пластинами радіуса $R$ заряджається струмом $I$ від зовнішнього
	джерела (рис.~\ref{pic:CapacitorPoynting}). Якщо в~деякий момент заряд на нижній
	пластині дорівнює $q$, то вектор індукції всередині конденсатора
	% .........................................................
	\begin{equation}\label{eq:CapacitorD}
		D = 4\pi\sigma = \frac{4\pi q}{\pi R^2},
	\end{equation}
	% .........................................................
	напрямлений від нижньої пластини до верхньої. Густина струму зміщення в~конденсаторі:
	% .........................................................
	\begin{equation*}
		j_\text{зм} = \frac{\dot D}{4\pi} = \frac{\dot q}{\pi R^2}.
	\end{equation*}
	% .........................................................
	Цей струм породжує в~конденсаторі магнітне поле, яке (як і~в~попередньому прикладі)
	знаходимо за теоремою про циркуляцію, взявши за контур $L$ коло радіуса $r$ з~центром
	на осі конденсатора:
	% .........................................................
	\begin{equation*}
		2\pi r H(r) = \frac{4\pi}{c}\pi r^2 j_\text{зм}(r) \quad\Rightarrow\quad H(r) = \frac{2r}{cR^2}\dot q.
	\end{equation*}
	% .........................................................
	Оскільки конденсатор заряджається ($\dot q>0$), струм зміщення напрямлений від нижньої обкладки до верхньої, а~силові лінії $\Hfield$ охоплюють вісь конденсатора за правилом правого гвинта.

	Вектор Пойнтінга $\vect\Pi=\dfrac{c}{4\pi}\left[\Efield\times\Hfield\right]$
	спрямований \emph{до осі} конденсатора (перпендикулярно до бічної поверхні,
	всередину) і~дорівнює за величиною на межі ($r=R$):
	% .........................................................
	\begin{equation*}
		\Pi = \frac{c}{4\pi}EH = \frac{c}{4\pi}\left(\frac{4\pi q}{\epsilon\pi R^2}\right)\left(\frac{2\dot q}{cR}\right) = \frac{2q\dot q}{\epsilon\pi R^3}.
	\end{equation*}
	% .........................................................
	Зарядку вважаємо повільною, а~зазор малим ($h\ll R$), тому магнітною енергією між
	обкладками й~крайовими ефектами нехтуємо. Якщо відстань між пластинами дорівнює $h$,
	енергія, що втікає в~конденсатор за одиницю часу (крізь бічну поверхню площею $2\pi
	Rh$):
	% .........................................................
	\begin{equation*}
		\frac{dW}{dt} = \Pi\cdot 2\pi Rh = \frac{2q\dot q}{\epsilon\pi R^3}\cdot 2\pi Rh = \frac{4hq\dot q}{\epsilon R^2} = \frac{q\dot q}{C},
	\end{equation*}
	% .........................................................
	де $C=\epsilon\pi R^2/(4\pi h)$~--- ємність конденсатора (розділ~\ref{ElEnergy}). За
	кінцевий час, коли заряд на пластині досягає значення $q$, енергія, що втекла
	в~конденсатор ззовні, становить
	% .........................................................
	\begin{equation}\label{eq:CapacitorEnergyResult}
		W = \frac{q^2}{2C}
	\end{equation}
	% .........................................................
	--- ми прийшли, зовсім іншим шляхом, до вже відомої формули для енергії зарядженого
	конденсатора~\eqref{eq:CapacitorEnergy}!
\end{Example}




\begin{SCfigure}[2][htbp]\centering
	\localinput{wire-poynting.tikz}
	\caption{Поля $\Efield$, $\Hfield$ і~вектор Пойнтінга $\vect\Pi$ на поверхні довгого провідника зі струмом\label{pic:WirePoynting}}
\end{SCfigure}

\begin{Example}{Довгий провідник}
	Розглянемо довгий прямий дріт радіуса $R$ і~довжини $\ell$, яким тече струм $I$ (рис.~\ref{pic:WirePoynting}).
	На зовнішній поверхні проводу присутнє магнітне поле $H = 2I/(cR)$
	(розділ~\ref{MagFieldVacuum}). Струм у~проводі виникає завдяки напрузі $U$ на його
	кінцях, що створює всередині (і, за неперервністю тангенціальної складової, на
	поверхні) поле $E = U/\ell$.

	Вектор Пойнтінга виявляється спрямованим \emph{до осі} провідника й~дорівнює за
	величиною
	% .........................................................
	\begin{equation*}
		\Pi = \frac{c}{4\pi}EH = \frac{c}{4\pi}\cdot\frac{U}{\ell}\cdot\frac{2I}{cR} = \frac1{2\pi\ell R}\,UI.
	\end{equation*}
	% .........................................................
	Оскільки площа бічної поверхні провідника дорівнює $S=2\pi\ell R$, повний потік
	енергії, що втікає в~провідник, становить $\Pi S = UI$~--- це в точності та сама
	потужність, що йде на джоулеві втрати в~проводі: $\Pi S = \mathcal P = UI$
	(розділ~\ref{SteadyCurrent}).
\end{Example}



Обидва приклади дають одну картину, контрінтуїтивну, але фундаментальну: \emph{енергія надходить до споживача не вздовж проводів, а~ззовні, крізь бічну поверхню}.

У~конденсаторі енергія $q^2/2C$, що запасається в~полі, втікає крізь бічну поверхню, а~не вздовж проводів чи силових ліній $\Efield$. Величина збігається з~розрахунком через роботу джерела (розділ~\ref{ElEnergy}), але картина \emph{шляху} енергії інша.

У~провіднику енергія для подолання опору приходить з~поля навколо дроту й~входить у~нього збоку. Дріт~--- не \enquote{труба}, якою тече енергія, а~лише напрямна для цього потоку. Уздовж дроту енергію несе поле \emph{зовні}: нормальна складова $\Efield$ від поверхневих зарядів разом з~$\Hfield$ дає поздовжню складову $\vect\Pi$.


%% ========================================================
\section{Скін-ефект}\hypertarget{skin-effect}{}
%% ========================================================


До цього моменту рівняння Максвелла поставали перед нами або як \emph{підсумок} уже
відомого (стаціонарні випадки~--- електро- і~магнітостатика), або як \emph{джерело}
якісних висновків (струм зміщення, симетрія, теорема Пойнтінга). Тепер подивимося на них
як на \emph{робочий інструмент}, узявши конкретну фізичну задачу: поведінку
електромагнітного поля у~провіднику за швидкозмінних процесів, на прикладі скін-ефекту.


% ---------------------------------------------------------
\subsection*{Означення та наслідки}
% ---------------------------------------------------------


\emphx{Скін-ефектом}\index{Скін-ефект} (від англ. \emph{skin}~--- шкіра) називають протікання струмів
високої частоти переважно в~тонкому поверхневому шарі провідника
(рис.~\ref{pic:SkinIllustration}).

\begin{figure}[htbp]\centering
	\localinput{skin.tikz}
	\caption{Розподіл густини струму в~перерізі мідного дроту діаметром $0{,}5$~мм на
	частотах $50$~Гц, $200$~кГц, $3$~МГц і~$100$~МГц (інтенсивність кольору пропорційна
	$|\vect j|$): зі зростанням частоти струмонесучий шар стоншується\label{pic:SkinIllustration}}
\end{figure}

Скін-ефект має кілька важливих практичних наслідків:
\begin{itemize}
	\item зменшується \emph{діюча} площа перерізу провідника, а~отже, \emph{зростає}
	  його опір на змінному струмі;
	\item \emph{зменшується} індуктивність провідника: оскільки за сильного скін-ефекту
	  струму (а~отже, і~магнітного поля) всередині провідника практично немає, магнітна
	  енергія стає меншою на величину енергії поля, яка раніше запасалась усередині
	  провідника;
	\item електромагнітна хвиля, що потрапляє на поверхню провідника (металу, електроліту,
	  плазми), швидко загасає вглиб нього, проникаючи лише на невелику глибину~--- це
	  тісно пов'язано з~тим, що \emph{провідники добре відбивають світло}.
\end{itemize}

\begin{Historical}
	Скін-ефект першим теоретично описав Горацій Ламб (1883)~--- для провідних куль; Олівер
	Хевісайд, який ще з~початку 1880-х років вивчав поширення сигналів у~телеграфних
	лініях, розглянув його для проводів (1885~р.). Отже, явище теоретично описали ще до
	масового поширення змінного струму в~енергетиці~--- у~зв'язку з~передаванням сигналів
	телеграфними лініями.
\end{Historical}


% ---------------------------------------------------------
\subsection*{Якісне пояснення}
% ---------------------------------------------------------


Змінний струм у~провіднику породжує змінне вихрове магнітне поле, силові лінії якого
перпендикулярні до осі провідника (рис.~\ref{pic:SkinEffectQualitative}). Це змінне
магнітне поле, своєю чергою, породжує вихрове електричне поле (розділ~\ref{ElMagInduction}),
що спричиняє в~товщі провідника струми Фуко: на \emph{поверхні} провідника вихровий струм
напрямлений \emph{у~той самий бік}, що й~основний струм провідності (і~\emph{підсилює}
його там), а~\emph{всередині} провідника~--- у~\emph{протилежний} бік
(і~\emph{послаблює} струм у~серцевині). У~результаті струм \emph{витісняється} з~осі
провідника до його поверхні.

\begin{SCfigure}[3][htbp]\centering
	\localinput{skin-effect.tikz}
	\caption{Якісний механізм скін-ефекту: вихрові струми Фуко підсилюють струм на
	поверхні й послаблюють у~центрі провідника\label{pic:SkinEffectQualitative}}
\end{SCfigure}


% ---------------------------------------------------------
\subsection*{Чому в~металах нехтують струмами зміщення}
% ---------------------------------------------------------


У~провіднику густина струму провідності й~напруженість поля пов'язані законом Ома, $\vect
j = \lambda\Efield$. Якщо поле в~провіднику змінюється гармонічно,
$\Efield=\Efield_0\cos(\omega t)$, то поряд зі струмом провідності виникає й~струм
зміщення:
% .........................................................
\begin{equation*}
	\vect j_\text{зм} = \frac{\epsilon}{4\pi}\parttime{\Efield} = \frac{\epsilon\omega}{4\pi}E_0\cos\left(\omega t+\frac\pi2\right).
\end{equation*}
% .........................................................
Оцінимо відношення амплітуд густини струму зміщення та струму провідності:
% .........................................................
\begin{equation}\label{eq:SkinDisplacementNegligible}
	\frac{j_\text{зм}}{j_\text{пров}} = \frac{\epsilon\omega}{4\pi\lambda} = \omega\tau_\text{М} \ll 1,
\end{equation}
% .........................................................
де $\tau_\text{М} = \dfrac{\epsilon}{4\pi\lambda}$~--- максвелівський час релаксації
\eqref{eq:MaxwellRelaxation}. Умова означає, що $\tau_\text{М} \ll T$, де $T \sim
1/\omega$~--- період зміни електромагнітного поля.

Для типових металів провідність $\lambda \sim 10^{16}$--$10^{18}$~с$^{-1}$ (гауссова
система), а~типова діелектрична проникність остова~--- порядку $10$--$100$, що відповідає
вкрай малому максвелівському часу $\tau_\text{М} \sim 10^{-19}$--$10^{-18}$~с. Тому формально нерівність~\eqref{eq:SkinDisplacementNegligible} виконується аж до $\omega\sim1/\tau_\text{М}\sim10^{18}$~с$^{-1}$; реальну межу задає не вона, а~умова $\omega\tau\ll1$ застосовності закону Ома (див. рамку нижче), тобто $\nu\lesssim10^{13}$--$10^{14}$~Гц. У~всьому цьому діапазоні струми зміщення \emph{не відіграють істотної ролі} у~металах порівняно зі струмами провідності~--- саме тому надалі, аналізуючи скін-ефект у~металах, ми зможемо ними знехтувати.

Варто уточнити фізичний зміст діелектричної проникності $\epsilon$ у~цій оцінці:
у~металі вона описує \emph{не} відгук вільних електронів провідності (їхній внесок
уже повністю врахований окремо, через закон Ома і~струм провідності), а~значно слабшу
поляризацію \emph{зв'язаних} електронів іонного остова. Крім того, сама провідність
$\lambda$ у~реальності не є сталою величиною~--- вона починає спадати з~частотою,
щойно період поля стає порівнянним із часом релаксації електронів провідності ($\tau
\sim 10^{-14}$~с). Слід розрізняти макроскопічний час $\tau_\text{М}$ та
мікроскопічний час вільного пробігу електронів $\tau$: оскільки формально розрахований
$\tau_\text{М} \ll \tau$, реальна швидкість релаксації в~металі обмежується інерцією
самих носіїв (моделлю Друде). Тому межа $\nu\sim10^{13}$~Гц~--- лише орієнтовна, за порядком величини, а~не різка межа. Формули цього підрозділу (\emph{нормальний} скін-ефект)
застосовні, поки $\omega\tau\ll1$; на вищих частотах доводиться враховувати частотну
залежність $\lambda(\omega)$.


% ---------------------------------------------------------
\subsection*{Рівняння скін-ефекту}
% ---------------------------------------------------------


Виведемо\index{Скін-ефект!глибина проникнення} рівняння для поля в~провіднику, спираючись безпосередньо на рівняння Максвелла (а~не на якісні міркування). Нехай $\epsilon=\const$, $\mu=\const$; у~провіднику
немає некомпенсованого об'ємного заряду, тому теорема Гаусса дає $\divg\Efield=0$.
З~урахуванням закону Ома, а~також знехтувавши, за щойно встановленою
нерівністю~\eqref{eq:SkinDisplacementNegligible}, струмом зміщення, рівняння Максвелла
в~металі набувають вигляду
% .........................................................
\begin{equation*}
	\rot\Efield = -\frac{\mu}c\parttime{\Hfield}, \qquad \rot\Hfield = \frac{4\pi\lambda}{c}\Efield, \qquad \divg\Efield=0, \quad \divg\Hfield=0.
\end{equation*}
% .........................................................
Виключимо $\Hfield$: застосуємо операцію $\rot$ до першого рівняння й~підставимо друге:
% .........................................................
\begin{equation*}
	\rot\rot\Efield = -\frac{\mu}c\parttime{}\left(\rot\Hfield\right) = -\frac{4\pi\mu\lambda}{c^2}\parttime{\Efield}.
\end{equation*}
% .........................................................
З~іншого боку, за тотожністю векторного аналізу і~з~урахуванням $\divg\Efield=0$:
$\rot\rot\Efield = \grad(\divg\Efield)-\nabla^2\Efield = -\nabla^2\Efield$. Прирівнюючи,
отримуємо \emphx{рівняння скін-ефекту}:
% .........................................................
\begin{equation}\label{eq:SkinWaveEquation}
	\tcbhighmath{
		\nabla^2\Efield = \frac{4\pi\mu\lambda}{c^2}\parttime{\Efield}.
	}
\end{equation}
% .........................................................
Таке саме рівняння (з~тими самими кроками виведення) отримуємо й~для $\Hfield$ та для густини струму $\vect j=\lambda\Efield$. Зверніть увагу: у~ньому \emph{перша} похідна за часом, а~не друга, як у~хвильовому рівнянні. Отже, поле в~металі не \enquote{біжить}, а~\emph{просочується} вглиб, як тепло.


% ---------------------------------------------------------
\subsection*{Оцінка глибини проникнення}
% ---------------------------------------------------------


Рівняння~\eqref{eq:SkinWaveEquation} за математичною формою збігається з~рівнянням
\emphx{дифузії}, $\parttime{F} = D\nabla^2 F$, з~коефіцієнтом дифузії $D =
c^2/(4\pi\mu\lambda)$; зокрема, з~рівнянням теплопровідності, де роль $F$ відіграє температура.

\emph{Аналогія.} Узимку ґрунт промерзає не одразу: добові коливання температури повітря проникають у~землю лише на ${\sim}10$~см, а~річні (сезонні)~--- на кілька метрів. Глибина проникнення зростає як $\sqrt{\chi T}$ ($\chi$~--- температуропровідність, $T$~--- період коливань), тому в~льоху на глибині кількох метрів майже постійна температура. Скін-шар~--- це такий самий \enquote{льох} для змінного поля: чим швидше змінюється поле на поверхні, тим тонший шар, до якого воно встигає \enquote{просочитися}.

Якщо на поверхні речовини діє джерело, то за час $\tau$
\enquote{хвиля дифузії} проникає на глибину $\delta=\sqrt{D\tau}$; якщо джерело діє
періодично з~періодом $T=2\pi/\omega$, слід узяти $\tau\sim1/\omega$, звідки отримуємо \emph{оцінку} для глибини
проникнення:
% .........................................................
\begin{equation}\label{eq:SkinDepthEstimate}
	\delta \sim \frac{c}{\sqrt{4\pi\lambda\mu\omega}}.
\end{equation}
% .........................................................

Переконаємось, що ця оцінка, з~точністю до числового множника, дійсно правильна,
застосувавши безпосередньо рівняння Максвелла в~інтегральній формі до тонкого
приповерхневого шару завтовшки $\delta$.

\emphx{Циркуляція $\Efield$.} Розглянемо прямокутний контур $L$ зі сторонами $h$ (уздовж
осі провідника, на його поверхні) і~$\delta$ (углиб) (рис.~\ref{pic:SkinEstimate1}). За
законом електромагнітної індукції, $Eh = -\dfrac1c\dfrac{\Delta\Phi}{\Delta t}$, де
магнітний потік крізь цей контур за модулем $\Phi\approx BS\approx\mu H h\delta$ (поле $H$ спадає до нуля на
глибині $\delta$), а~швидкість його зміни за порядком величини $\Delta\Phi/\Delta t
\approx \omega\Phi$ ($\omega\sim1/T$~--- частота зміни струму). Тому
% .........................................................
\begin{equation}\label{eq:SkinCirculationE}
	Eh \sim \frac1c\mu H\omega h\delta \quad\Rightarrow\quad E \sim \frac{\mu H\omega\delta}{c}.
\end{equation}
% .........................................................

\emphx{Циркуляція $\Hfield$.} Тепер розглянемо контур $L$ на поперечному перерізі
провідника, завдовжки $\ell$ уздовж поверхні й~завширшки $\delta$ углиб
(рис.~\ref{pic:SkinEstimate2}). За теоремою про циркуляцію, $H\ell =
\dfrac{4\pi}{c}j\ell\delta$, а~з~огляду на закон Ома, $j=\lambda E$:
% .........................................................
\begin{equation}\label{eq:SkinCirculationH}
	H \sim \frac{4\pi\lambda}{c}E\delta.
\end{equation}
% .........................................................

\begin{figure}[htbp]
	\centering
	\captionbox{Контур для оцінки циркуляції $\Efield$ поблизу поверхні циліндричного провідника (вигляд збоку)\label{pic:SkinEstimate1}}[0.47\textwidth]{%
		\localinput{skin-estimate-1.tikz}%
	}\hfill
	\captionbox{Контур для оцінки циркуляції $\Hfield$ поблизу поверхні циліндричного провідника (вигляд зверху)\label{pic:SkinEstimate2}}[0.47\textwidth]{%
		\localinput{skin-estimate-2.tikz}%
	}
\end{figure}

Підставляючи~\eqref{eq:SkinCirculationH} в~\eqref{eq:SkinCirculationE}:
% .........................................................
\begin{equation*}
	E \sim \frac{4\pi\mu\lambda\omega\delta^2}{c^2}\, E,
\end{equation*}
% .........................................................
звідки, скорочуючи на $E$, отримуємо оцінку глибини скін-шару:
% .........................................................
\begin{equation}\label{eq:SkinDepthFinal}
	\delta \sim \frac{c}{\sqrt{4\pi\lambda\mu\omega}},
\end{equation}
% .........................................................
що збігається з~дифузійною оцінкою~\eqref{eq:SkinDepthEstimate}. Точне значення числового коефіцієнта дає розв'язок рівняння~\eqref{eq:SkinWaveEquation} для однорідного провідника, що займає півпростір $x>0$. Шукаючи його у~вигляді плоскої хвилі $\Efield\propto e^{i(kx-\omega t)}$, отримуємо $k^2=4\pi i\mu\lambda\omega/c^2$, тобто $k=(1+i)/\delta$, де
% .........................................................
\begin{equation}\label{eq:SkinDepthExact}
	\tcbhighmath{
		\delta = \frac{c}{\sqrt{2\pi\lambda\mu\omega}}.
	}
\end{equation}
% .........................................................
Поле загасає вглиб провідника за законом $E\propto e^{-x/\delta}\cos(x/\delta-\omega t)$, а~$\delta$~--- глибина, на якій амплітуда зменшується в~$e$ разів. Оцінка~\eqref{eq:SkinDepthFinal} відрізняється від точної формули лише множником $\sqrt2$.
Значення $\delta$ для деяких металів на частоті $10$~ГГц, обчислені за точною формулою,
наведено в~табл.~\ref{tab:SkinDepth}.

\begin{Attention}[unbreakable]
	Оцінка~\eqref{eq:SkinDepthFinal} побудована на спрощених контурних міркуваннях і~дає
	лише порядок величини та правильну залежність
	$\delta\propto1/\sqrt{\lambda\mu\omega}$; точний числовий коефіцієнт дає розв'язок
	рівняння~\eqref{eq:SkinWaveEquation}. Формула~\eqref{eq:SkinDepthExact}
	описує \emph{нормальний} скін-ефект: вона застосовна, коли $\delta$ значно менша за
	розміри й~радіус кривини провідника, частота така, що $\omega\tau_0\ll1$ ($\tau_0\sim10^{-14}$~с~--- час вільного пробігу електрона), і~довжина
	вільного пробігу електрона менша за $\delta$ (тоді справджується локальний закон Ома
	$\vect j=\lambda\Efield$). Інакше (наприклад, у~чистих металах при низьких
	температурах) спостерігають \emph{аномальний} скін-ефект.
\end{Attention}

\begin{table}[htbp]\small\centering
	\caption{Глибина скін-шару для деяких металів на частоті $10$~ГГц (за точною
	формулою)\label{tab:SkinDepth}}
	\begin{tblr}{
			colspec={X[l, m]X[c,m]},
			row{odd} = {gray!10},
			row{1} = {c, m, fg=white, bg=themecolorlight, font=\bfseries},
			hlines,
			vlines
		}
		Метал              & Глибина скін-шару, мкм \\
		Мідь (\ce{Cu})     & 0{,}652                \\
		Срібло (\ce{Ag})   & 0{,}634                \\
		Золото (\ce{Au})   & 0{,}753                \\
		Алюміній (\ce{Al}) & 0{,}820                \\
	\end{tblr}
\end{table}

\emph{Числові орієнтири.} Для міді $\delta\approx66~\text{мм}/\sqrt{f[\text{Гц}]}$: на частоті мережі $50$~Гц $\delta\approx9$~мм, на $1$~МГц~--- $66$~мкм, на $10$~ГГц~--- менше мікрометра. Тому для звичайної проводки скін-ефект неістотний, а~от для шин електростанцій і~проводів ЛЕП завтовшки кілька сантиметрів~--- уже помітний. У~сталі через велике $\mu\sim10^3$ скін-шар у~десяток і~більше разів тонший, ніж у~міді: на частоті індукційної плити ($20$--$100$~кГц) він становить лише десятки мікрометрів, тож дно сталевої каструлі нагрівається тонким поверхневим шаром, а~ось алюмінієва чи мідна каструля на такій плиті майже не гріється.


% ---------------------------------------------------------
\subsection*{Застосування скін-ефекту в~техніці}
% ---------------------------------------------------------


Існування\index{Скін-ефект!застосування} скін-ефекту завжди враховують у~техніці швидкозмінних струмів. Оскільки такі
струми практично не йдуть у~глибині провідника, проводи й~шини для них \emph{виготовляють
із порожнистих труб}~--- економлячи метал без істотної втрати в~опорі. У~радіотехніці
надвисоких частот багато деталей (хвилеводи, коаксіальні лінії) покривають тонким, добре
провідним шаром срібла, оскільки їхній опір практично повністю визначається лише
поверхневим шаром. Скін-ефект у~поєднанні з~нагріванням вихровими струмами
(розділ~\ref{ElMagInduction}) використовують і~для \emph{поверхневого загартування сталі}:
деталі піддають дії змінних струмів і~потім швидко охолоджують, унаслідок чого розігрів
і~загартування відбуваються лише в~поверхневому шарі, а~основна маса деталі лишається
незагартованою~--- це дає вироби з~високою стійкістю до стирання поверхні, але без
крихкості, властивої звичайному наскрізному загартуванню. Щоб, навпаки, \emph{послабити}
скін-ефект у~котушках високої частоти, провід роблять із багатьох тонких ізольованих
жилок, сплетених так, щоб кожна по черзі опинялася на поверхні джгута (\emph{літцендрат}).


%% ========================================================
\section*{Підсумки}
%% ========================================================


\phantomsection \addcontentsline{toc}{section}{Підсумки}
%% --------------------------------------------------------

У~цьому розділі ідея Максвелла про струм зміщення замкнула систему рівнянь електродинаміки
в~остаточному вигляді й~дала нам два перші, найближчі до вершини, наслідки цієї системи:
закон збереження енергії електромагнітного поля у~формі теореми Пойнтінга та кількісний
опис скін-ефекту у~провідниках. В~обох випадках рівняння Максвелла виступили вже не
підсумком відомого, а~робочим інструментом, що дає нові, нетривіальні результати~--- аж до
контрінтуїтивного висновку про те, що енергія до провідника надходить не вздовж нього,
а~ззовні, крізь бічну поверхню.

\textbf{Що запам'ятати.}
\begin{itemize}[nosep,leftmargin=1.2em]
	\item Змінне електричне поле породжує магнітне так само, як струм: $\vect j_\text{зм}=\dfrac1{4\pi}\parttime{\Dfield}$. Повний струм $\vect j+\vect j_\text{зм}$ завжди замкнений.
	\item Струм зміщення не переносить заряд і~не виділяє джоулевого тепла; у~металах аж до інфрачервоних частот ним можна знехтувати.
	\item Енергія поля має густину $w=(\Efield\Dfield+\Bfield\Hfield)/8\pi$ і~тече з~густиною потоку $\vect\Pi=\frac{c}{4\pi}[\Efield\times\Hfield]$; у~провідник вона входить \emph{збоку}.
	\item Змінне поле \enquote{просочується} в~метал дифузійно, на глибину $\delta=c/\sqrt{2\pi\lambda\mu\omega}\propto1/\sqrt\omega$; для міді на $50$~Гц це близько $1$~см.
\end{itemize}

Наступний крок~--- спуститися з~вершини далі, у~простір хвильових явищ: розглянути, як із
рівнянь Максвелла у~вакуумі випливає саме існування електромагнітних хвиль, що біжать зі
швидкістю світла, а~тоді~--- дослідити, як саме заряджені частинки, що рухаються
з~прискоренням, ці хвилі випромінюють.


%% ========================================================
\section*{Задачі}
%% ========================================================


\phantomsection \addcontentsline{toc}{section}{Задачі}
%% --------------------------------------------------------

% задачі цього розділу короткі: не розриваємо їх між сторінками
\begingroup
\tcbset{answersolutionstyle/.append style={unbreakable}}

\begin{problem}\label{prb:CapacitorDisplacementAC}
	\emph{Струм зміщення в~конденсаторі змінної напруги.} Плоский конденсатор із круглими
	пластинами радіуса $R$, відстанню між ними $d$ і діелектриком проникності $\epsilon$
	приєднано до джерела напруги $U(t)=U_m\sin\omega t$. Знайдіть густину й~повний струм
	зміщення між пластинами та напруженість магнітного поля $H$ на краю пластин ($r=R$).
	Порівняйте струм зміщення зі струмом $C\,dU/dt$ у~підвідних дротах ($C=\epsilon
	R^2/4d$~--- ємність конденсатора). \emph{Вказівка:} $E=U/d$; далі скористайтесь
	означенням $j_\text{зм}$ та теоремою про циркуляцію.
	\begin{solution}
		$j_\text{зм}=\dfrac{\epsilon\omega U_m}{4\pi d}\cos\omega t$,
		$I_\text{зм}=\dfrac{\epsilon R^2\omega U_m}{4d}\cos\omega t = C\dfrac{dU}{dt}$ (тобто дорівнює струму в дротах --- струм зміщення \enquote{замикає} коло);
		$H(R)=\dfrac{2I_\text{зм}}{cR}=\dfrac{\epsilon R\omega U_m}{2cd}\cos\omega t$.
	\end{solution}
\end{problem}

\begin{problem}\label{prb:MetalDisplacementFreq}
	\emph{Коли струм зміщення в~металі стає помітним?} Оцініть частоту $\nu$, на якій амплітуди густин струму зміщення й~струму провідності в~міді зрівнюються ($\lambda\approx5\cdot10^{17}$~с$^{-1}$, $\epsilon\sim10$). Порівняйте з~межею застосовності закону Ома, $\omega\tau\sim1$ ($\tau\sim10^{-14}$~с), і~зробіть висновок. \emph{Вказівка:} прирівняйте $\epsilon\omega/4\pi$ до $\lambda$.
	\begin{solution}
		$\omega\sim4\pi\lambda/\epsilon\approx6\cdot10^{17}~\text{с}^{-1}$, тобто $\nu=\omega/2\pi\approx10^{17}$~Гц~--- далекий ультрафіолет чи м'який рентген. Закон Ома ж порушується вже за $\omega\sim1/\tau\sim10^{14}~\text{с}^{-1}$ ($\nu\sim10^{13}$~Гц). Отже, поки справджується $\vect j=\lambda\Efield$, струм зміщення в~металі нехтовно малий (оцінка орієнтовна, бо $\lambda$, $\epsilon$ і~$\tau$ відомі лише за порядком величини).
	\end{solution}
\end{problem}

\begin{problem}\label{prb:SolenoidPoynting}
	\emph{Енергія втікає в~соленоїд.} По довгому соленоїду радіуса $R$ з~$n$ витками на
	одиницю довжини тече струм $I(t)$, що повільно зростає. Знайдіть вектор Пойнтінга на
	бічній поверхні соленоїда й~покажіть, що потік енергії крізь цю поверхню дорівнює
	швидкості зростання магнітної енергії всередині~--- так само як у~прикладі
	з~конденсатором. \emph{Вказівка:} $H=4\pi nI/c$ усередині; електричне поле на
	поверхні знайдіть із закону індукції Фарадея (циркуляція $\Efield$ по колу радіуса
	$R$).
	\begin{solution}
		$E=\dfrac{2\pi nR\dot I}{c^2}$
		(вздовж кола), $\Pi=\dfrac{2\pi n^2R\,I\dot I}{c^2}$, вектор $\vect\Pi$ спрямований \emph{до осі}; на одиницю довжини соленоїда
		$\Pi\cdot2\pi R=\dfrac{4\pi^2n^2R^2}{c^2}I\dot I=\dfrac{d}{dt}\!\left(\dfrac{L_1I^2}{2c^2}\right)$, де $L_1=4\pi^2n^2R^2$ --- індуктивність
		одиниці довжини.
	\end{solution}
\end{problem}

\begin{problem}\label{prb:SkinDepthCopper50Hz}
	\emph{Скін-шар на промисловій частоті.} Знаючи з~табл.~\ref{tab:SkinDepth} глибину скін-шару міді на частоті $10$~ГГц
	($0{,}652$~мкм), знайдіть її на частоті $50$~Гц і~зробіть висновок про доцільність суцільних мідних проводів великого діаметра
	в~лініях електропередачі. \emph{Вказівка:} $\delta\propto1/\sqrt\omega$.
	\begin{solution}
		$\delta_{50}=\delta_{10\ \text{ГГц}}\sqrt{10^{10}/50}\approx 0{,}652~\text{мкм}\cdot1{,}4\cdot10^4\approx 9$~мм; отже, струм майже не йде глибше ніж на $\sim1$~см, і~суцільні мідні (чи алюмінієві) дроти діаметром понад $\sim2$~см малоефективні~--- тому їх роблять порожнистими або складеними з~тонших жил.
	\end{solution}
\end{problem}

\begin{problem}\label{prb:CoaxPoynting}
	\emph{Енергія в~коаксіальному кабелі.} Між внутрішнім (радіуса $a$) і~зовнішнім (радіуса $b$) провідниками ідеального коаксіального кабелю прикладено напругу $U$, а~по ньому тече струм $I$. Знайдіть вектор Пойнтінга між провідниками й~покажіть, що потік енергії крізь переріз кабелю дорівнює $UI$. \emph{Вказівка:} $E=U/(r\ln(b/a))$, $H=2I/(cr)$.
	\begin{solution}
		$\Pi=\dfrac{c}{4\pi}EH=\dfrac{UI}{2\pi r^2\ln(b/a)}$ вздовж кабелю; $\displaystyle\int_a^b\Pi\,2\pi r\,dr=\dfrac{UI}{\ln(b/a)}\int_a^b\dfrac{dr}{r}=UI$~--- енергія тече не в~провідниках, а~в~проміжку між ними.
	\end{solution}
\end{problem}
\endgroup

\begin{Questions}
	\begin{quizlist}
		\item У~чому полягала ідея Максвелла про симетрію електричного й~магнітного полів?
		  Як ця ідея привела його до поняття струму зміщення?
		\item Покажіть, що теорема про циркуляцію $\Hfield$ для постійних полів,
		  застосована до випадку змінного в~часі заряду, суперечить закону збереження
		  заряду (рівнянню неперервності).
		\item Виведіть вираз для густини струму зміщення $\vect j_\text{зм} =
		  \dfrac1{4\pi}\parttime{\Dfield}$ з~вимоги узгодженості теореми про циркуляцію із
		  законом збереження заряду. Чому цей вибір не єдиний математично можливий, і~що
		  зрештою підтверджує саме його правильність?
		\item Запишіть остаточну теорему про циркуляцію магнітного поля з~урахуванням
		  струму зміщення в~інтегральній та диференціальній формах.
		\item Розгляньте задачу про радіальне стікання заряду з~кулі: чому магнітне поле
		  в~цій задачі не виникає взагалі, попри те що заряд явно рухається? Яку роль тут
		  відіграє точна компенсація струму провідності струмом зміщення?
		\item Розгляньте задачу про зарядку плоского конденсатора: чому циркуляція
		  $\Hfield$ по одному й~тому самому контуру дає однаковий результат для двох
		  різних поверхонь~--- тієї, що перетинає провід, і~тієї, що проходить між
		  обкладками? Що це означає для ролі струму зміщення в~конденсаторі?
		\item Чим струм зміщення принципово відрізняється від струму провідності~--- чому
		  до нього незастосовні закон Ома і~закон Джоуля--Ленца?
		\item Порівняйте закон електромагнітної індукції Фарадея та гіпотезу Максвелла про
		  магнітоелектричну індукцію. У~чому полягає асиметрія знаків між ними та до яких
		  глибоких наслідків вона приводить?
		\item Запишіть повну систему чотирьох рівнянь Максвелла в~інтегральній та
		  диференціальній формах. За яким принципом ці рівняння групуються в~дві пари?
		\item Покажіть, що закон збереження заряду випливає з~рівнянь Максвелла
		  (з~\eqref{eq:MaxwellGaussE} і~\eqref{eq:MaxwellAmpere}).
		\item Виведіть вираз електричного поля через потенціали, $\Efield = -\grad\phi -
		  \dfrac1c\parttime{\vect A}$, з~теореми Гаусса для $\Bfield$ і~закону індукції
		  Фарадея. Чому в~нестаціонарному випадку $\Efield$ визначається не лише скалярним
		  потенціалом?
		\item Покажіть, що в~стаціонарному випадку повна система рівнянь Максвелла
		  розпадається на незалежні системи електростатики й~магнітостатики. Чому зв'язок
		  електрики й~магнетизму~--- принципово нестаціонарне явище?
		\item Чому потенціали $\phi$ і~$\vect A$ визначено неоднозначно? Як перетворюються
		  вони при зміні калібрування, щоб поля $\Efield$ і~$\Bfield$ не змінилися?
		\item Перелічіть матеріальні рівняння, що доповнюють рівняння Максвелла
		  ($\Dfield=\epsilon\Efield$, $\Bfield=\mu\Hfield$, закон Ома). Навіщо вони
		  потрібні для розв'язання конкретних задач?
		\item Виведіть граничні умови для електричних і~магнітних векторів на межі розділу двох середовищ (циліндр Гаусса, контур циркуляції). Чому змінність полів не змінює їхнього вигляду? Коли $\Hfield_\tau$ зазнає стрибка?
		\item Яка структура рівнянь Максвелла? Що породжує $\Dfield$, що~--- $\Hfield$, і~чому магнітне поле не має зарядів-джерел (рис.~\ref{pic:MaxwellScheme})?
		\item Виведіть теорему Пойнтінга $-\parttime{w} = \divg\vect\Pi + \vect
		  j\cdot\Efield$ з~рівнянь Максвелла для $\rot\Hfield$ і~$\rot\Efield$,
		  використовуючи тотожність для дивергенції векторного добутку.
		\item Що таке вектор Пойнтінга і~густина енергії електромагнітного поля?
		  Сформулюйте теорему Пойнтінга в~інтегральній формі та поясніть фізичний зміст
		  кожного доданка.
		\item Розгляньте приклад зарядки конденсатора через теорему Пойнтінга: чому вся
		  енергія, що запасається в~полі всередині конденсатора, насправді втікає туди
		  ззовні, крізь бічну поверхню, а~не вздовж проводів?
		\item Розгляньте приклад довгого провідника зі струмом: чому енергія, що йде на
		  джоулеві втрати, надходить у~провідник не вздовж нього від джерела, а~ззовні,
		  крізь бічну поверхню?
		\item Що таке скін-ефект і~які його основні практичні наслідки (зміна опору,
		  індуктивності, відбивання світла провідниками)?
		\item Дайте якісне пояснення скін-ефекту через вихрові струми Фуко: чому вони
		  підсилюють струм на поверхні провідника й~послаблюють у~його серцевині?
		\item Чому в~металах аж до дуже високих частот ($\sim10^{13}$ Гц) можна нехтувати
		  струмом зміщення порівняно зі струмом провідності? Виведіть умову
		  $\epsilon\omega/(4\pi) \ll \lambda$.
		\item Виведіть рівняння скін-ефекту $\nabla^2\Efield =
		  (4\pi\mu\lambda/c^2)\,\parttime{\Efield}$ з~рівнянь Максвелла для провідника,
		  знехтувавши струмом зміщення.
		\item Виведіть оцінку глибини скін-шару $\delta \sim c/\sqrt{4\pi\lambda\mu\omega}$, спираючись на аналогію з~рівнянням дифузії, а~також прямим контурним міркуванням через циркуляції $\Efield$ і~$\Hfield$ біля поверхні провідника.
		\item Які практичні застосування скін-ефекту в~техніці (порожнисті провідники ЛЕП, посріблення хвилеводів, поверхневе загартування сталі)?
	\end{quizlist}
\end{Questions}

\begin{Questions}[Питання на розуміння]
	\begin{quizlist}
		\item Між обкладками конденсатора, який заряджається, поставили компас. Чи відхилиться стрілка? Що зміниться, якщо зарядка припиниться?
		\item Батарейку з'єднали з~лампочкою двома довгими проводами. Покажіть на рисунку напрямок $\Efield$, $\Hfield$ і~$\vect\Pi$ поблизу проводів. Де саме тече енергія до лампочки?
		\item Чи можна вважати, що енергія в~колі постійного струму тече \emph{всередині} проводів, як вода в~трубі? Що змінилося б, якби поле зовні проводів було відсутнє?
		\item Чому на індукційній плиті добре гріється сталева каструля, але погано~--- алюмінієва, хоча алюміній проводить струм краще?
		\item Мікрохвильову піч зсередини можна було б зробити з~пластику, покритого шаром срібла завтовшки $1$~мкм. Чи екранував би такий шар випромінювання з~частотою $2{,}45$~ГГц? Оцініть.
		\item У~скільки разів зміниться глибина скін-шару, якщо частоту збільшити в~$100$ разів? Як це пов'язати з~тим, що добові коливання температури проникають у~ґрунт значно мілкіше, ніж річні?
	\end{quizlist}
\end{Questions}

\Closesolutionfile{answer}
